\documentclass[11pt]{article}
\usepackage{fontspec}
\usepackage[french]{babel}
\usepackage{amsmath}
\usepackage{amssymb}
\usepackage{amsthm}
\usepackage{mathrsfs}
\usepackage{tikz-cd}
\usepackage{enumitem}

\newcommand{\printedpage}[1]{\marginpar{\footnotesize p.~#1}}
\newcommand{\msp}[1]{\marginpar{\footnotesize [#1]}}
\newcommand{\Fl}{\mathrm{Fl}}
\newcommand{\Cat}{\mathrm{Cat}}
\newcommand{\card}{\operatorname{card}}

\frenchsetup{AutoSpacePunctuation=false}
\setlength{\parindent}{0pt}
\setlength{\parskip}{1ex}

\begin{document}

\printedpage{14}
% head: A. Grothendieck : Dérivateurs, Ch. XII, Édition des Universités de Montpellier II et Paris VII
attendais le moins : l'axiome des fibrations. Prenons la forme faible : Soit $f : X \longrightarrow Y$ une $W_\omega$-fibration et $f \in W_i$, les fibres $X_s$ sont elles $W_i$-asphériques (\footnote[37]{$\omega > \infty$, $W_\omega \subset W_\infty \subset W_n$})? On peut supposer $X$ et $Y$ 0-connexes, soit $x \in X_s$, la suite exacte d'homotopie donne
\[
\pi_{i+1}(X) \xrightarrow[\text{si } i \leq n-1]{\text{iso}} \pi_{i+1}(Y) \longrightarrow \pi_i(X_s) \longrightarrow \pi_i(X) \xrightarrow[\text{si } i \leq n]{\text{iso}} \pi_i(Y),
\]
donc
\[
\pi_i(X_s) = 1 \qquad \text{si } 0 \leq i \leq n-1,
\]
mais pour $i = n$ on trouve\msp{page 17}
\[
\pi_n(X_s) = \mathrm{Coker}\left(\pi_{n+1}(X) \longrightarrow \pi_{n+1}(Y)\right),
\]
il n'y a aucune raison que ce soit nul, comme ça devrait si $X_s \longrightarrow e$ était dans $W_n$.

Il est vrai que parfois on renforce justement la condition pour $W_n$, disons qu'on définit $W'_n \subseteq W_n$ par
\[
f \in W'_n \iff \left\{
\begin{array}{ll}
\text{a)} & \pi_0(f) \text{ bijectif,} \\[1ex]
\text{b)} & \pi_i(f,x) : \pi_i(X,x) \longrightarrow \pi_i(Y,f(x)) \text{ bijectif pour } 0 < i \leq n \text{ et } x \in X, \\[1ex]
\text{c)} & \pi_{n+1}(f,x) : \pi_{n+1}(X,x) \longrightarrow \pi_{n+1}(Y,f(x)) \ \textit{surjectif,}
\end{array}
\right.
\]
c'est à dire aussi
\[
f \in W'_n \iff \left\{
\begin{array}{l}
\text{les `fibres homotopiques' } Z \text{ de } X \longrightarrow Y \text{ (i.e. les fibres de } \bar{X} \xrightarrow{\bar{f}} Y\text{)} \\
\text{sont } n\text{-connexes, i.e. 0-connexe et } \pi_i(Z_y,z) = 0 \text{ pour } 1 \leq i \leq n.
\end{array}
\right.
\]
Alors $X$ dans Cat sera $W'_n$-asphérique si et seulement s'il est $W_n$-asphérique, i.e. $n$-connexe. Donc par la forme 2 de la définition de $W'_n$, l'axiome W(4b$''$) est satisfait. Il faut croire que ça pèche ailleurs! Peut-être déjà la saturation.

\hspace*{3em}Si $gf$ et $g$ sont dans $W'_i$, il n'est pas forcément de même de $f$,

car si $\bar{g}\bar{f}$ (et donc $\bar{g}$) est surjectif, il n'en est par forcément de même pour $\bar{f}$! Donc $W'_n$ ne mérite pas le nom de `localiseur'. % [sic] par forcément -> pas forcément, misprint

Par contre, $W_n$ me paraît excellent, si on ne lui demande pas l'impossible - de satisfaire à un `axiome de fibration' qui semble bien (à peu de choses près) caractériser le seul $W_\infty$ parmi les `bons' localiseurs. En tous cas l'axiome de connexité W(5) est satisfait, et l'axiome des limites filtrantes W(6) aussi. Par ailleurs W(1) (et même la saturation forte), W(2), W(3) sont satisfaits par tout $W_{\mathbf{D}}$ associé à un dérivateur - ce qui confirme à nouveau leur importance. Les axiomes W(4) (l'axiome `draconien' des fibrations, $\pm$ spécial à $W_\infty$), W(5) (connéxité) et % [sic] connéxité -> connexité, misspelling (written connexité elsewhere)
\printedpage{15}
% head: A. Grothendieck : Dérivateurs, Ch. XII, Édition des Universités de Montpellier II et Paris VII
\msp{page 18}W(6) (limites) m'apparaissent comme de nature plus spéciale, je ne songerais à imposer aucun d'eux pour l'usage terminologique de `localiseur fondamental'.

\newpage
\printedpage{16}
% head: A. Grothendieck : Dérivateurs, Ch. XII, Édition des Universités de Montpellier II et Paris VII
\section*{II\quad Perplexités sur les axiomes et la terminologie.}

Dans I et la plupart du temps, en travaillant avec les dérivateurs, j'ai supposé \fbox{$\ W_{\mathbf{D}} \supseteq W_\infty\ $}, hypothèse qui m'a parue essentielle (à tort (\footnote[38]{à \textit{tort} il s'avère maintenant.}) ou à raison). On voit donc que si $W_{\mathbf{D}}$ satisfait à l'axiome de fibrations W(4), même sous sa forme la plus faible W(4b$''$), il est \textit{égal} à $W_\infty$. Cela montre bien que cet axiome doit être vu comme pratiquement une \textit{caractérisation} de $W_\infty$ (parmi les localiseurs fondamentaux qui contiennent $W_\infty$ -- et ces localisateurs me paraissent les plus importants, et peut-être les seuls vraiment utiles).

La théorie des \textbf{D}-fibrations dans I, p. 111-126, a été faite dans l'esprit qu'on vient de dire, en supposant $W_{\mathbf{D}} \supseteq W_\infty$. La notion de \textbf{D}-fibrations à laquelle je m'étais finalement arreté était par suite différente des notions envisagées ici dans un esprit différent (page 4). Elle revient à tester la condition a) (trop forte presque toujours, sauf pour $W = W_\infty$), caractérisant les $W$-fibrations ultrafortes, par des $Y'' \longrightarrow Y'$, où $Y''$, $Y'$ sont tous deux $W_\infty$-asphériques (ce qui est plus exigeant que $W$-asphérique). C'est seulement quand $Y''$ et $Y'$ sont $W_\infty$-asphériques (au lieu de $W$-asphériques), qu'on demande que $X'' \longrightarrow X'$ soit dans $W$. On trouve donc\msp{page 19} une notion moins exigeante que celle de $W$-fibration forte. Si on dénote ici cette condition par b), on trouve (sous réserve que $W_\infty \subseteq W$), avec les notations de la page 4 (\footnote[39]{Cf. la bonne notion de $W$-fibration dans un contexte où on ne suppose pas $W_\infty \subseteq W$, dans th. 3 et déf. 3, p. 51, 52.}) % [sic] arreté -> arrêté, missing circumflex
% In the print, a curved arrow starts at "b" in the diagram below and points up to "b)" in the line "par b), on trouve ..." above; it is drawn here to the empty node at the top left.
\[
\begin{tikzcd}[column sep=small, row sep=large]
{} & & & & \\
& & \mathrm{a} \arrow[d, Rightarrow] & & W\text{-fibration ultraforte} \\
& & \mathrm{a}' \arrow[d, Rightarrow] & & W\text{-fibration forte} \\
& & \mathrm{b} \arrow[dl, Rightarrow] \arrow[dr, Rightarrow] \arrow[uuull, bend left=30] & & W\text{-fibration} \\
W\text{-fibration gauche} & \mathrm{a}_1 & & \mathrm{a}_2 & W\text{-fibration droite.}
\end{tikzcd}
\]
a$_1$), a$_2$) étaient aussi appelées \textbf{D}-fibrations cohomologique resp. \textbf{D}-fibrations homologiques (\footnote[40]{Non, c'est une notion encore plus forte, à coups de coefficients localement constants, cf. I, p. 121.}). Quand \textbf{D} satisfait Der 3a$'$ (I, p. 111) sur les \textbf{D}-coefficients localement constants sur les $X$ $W_\infty$-asphériques, alors les notions b), a$_1$), a$_2$) sont équivalentes. (Je présume que c'est le cas pour tous les \textbf{D} et tous les $W$ vraiment intéressants.) Par contre, j'ai l'impression que la notion de $W$-fibration forte, voire ultraforte, n'est guère utilisable que dans le cas $W = W_\infty$, où elle coïncide avec les précédentes. C'est dire qu'on peut probablement oublier la terminologie, et se rappeler par contre que dans le cas $W = W_\infty$, les % [sic] cohomologique -> cohomologiques, agreement with "D-fibrations"
\printedpage{17}
% head: A. Grothendieck : Dérivateurs, Ch. XII, Édition des Universités de Montpellier II et Paris VII
$W$-fibrations ont des propriétés d'une force extrême, s'exprimant par les deux propriétés dégagées dans W(4a) et W(4b).

\msp{page 20}Dans la théorie de $\mathrm{HOT}_W$, influencé par la réflexion générale sur les dérivateurs et l'axiome Der 3b, j'ai travaillé sous l'hypothèse $W \supseteq W_\infty$, mais en utilisant un axiome de fibration L5 $=$ W(4b) qui, il s'avère, implique finalement $W = W_\infty$. J'ai l'impression que je n'ai jusqu'à présent (propriétés de $f_!$ pour HOT) pas utilisé $W \supseteq W_\infty$, donc ça peut être utile, comme `ça ne coûte pas plus cher' (et oblige d'être plus attentif à ce qu'on utilise), de laisser tomber $W \supseteq W_\infty$ (impliquant $W = W_\infty$), et de mettre par contre sur $W$ des axiomes aussi forts (notamment de fibration) que les propriétés de $W_\infty$, `sans se priver en rien'. Je supposerai donc W(1), W(2), W(3), W(4a) (la forme la plus forte de l'axiome de fibration, cf. pages 13, 14) - c'est tout ce qui a été utile jusqu'à présent. Si on y joint W(5) (axiome de connexité) ça implique $W \subseteq W_\infty$. Mais pour l'instant je ne crois pas en avoir un besoin.

Mais méditant sur la question, je crois avoir trouvé le

\textbf{Théorème 2} (!) (\footnote[41]{Cf. pour un résultat plus fort XV 9.12, p. 133 bis. \textbf{NB} Le résultat est plus général, mais pas vraiment plus fort, car il ne donne pas $W = \Fl(\Cat)$.}) : \textit{Soit} $W \subseteq \Fl(\Cat)$ \textit{satisfaisant} W(1,2,3,4b$''$)\textit{. Si} $W$ \textit{ne satisfait pas à} W(5)\textit{, alors} $W = \Fl(\Cat)$.

\textsc{Démonstration}. Soit $f : X \longrightarrow Y$ dans Cat, tel que $\pi_0(f)$ ne soit pas bijectif. Soit $I$ la catégorie discrète ayant pour ensemble d'objets l'ensemble $\pi_0(Y)$, on a donc un diagramme commutatif
\[
\begin{tikzcd}[column sep=small]
X \arrow[rr, "f"] \arrow[dr, "p"'] & & Y \arrow[dl, "q"] \\
& I, &
\end{tikzcd}
\]
d'autre part un foncteur d'une catégorie dans une catégorie discrète est toujours $W_\omega$-parfait, donc $p$, $q$ sont des $W$-fibrations. Comme\msp{page 21} $f \in W$, il s'ensuit $f_i \in W$ pour tout $i \in I$ (\footnote[42]{On utilise W(4) sous forme plus forte que W(4b$''$)!}).

Si $\pi_0(f)$ n'était pas surjective, prenant $i$ qui n'est pas dans l'image, on aurait $X_i = \emptyset \longrightarrow Y_i \neq \emptyset$, qui serait dans $W$. Soit donc
\[
f : X = \emptyset \longrightarrow Y
\]
dans $W$. Alors $\forall$ $Z$ dans Cat, $f \times \mathrm{id}_Z : \underbrace{X \times Z}_{= \emptyset} \longrightarrow Y \times Z$ est dans $W$ (car $Z$ est
\printedpage{18}
% head: A. Grothendieck : Dérivateurs, Ch. XII, Édition des Universités de Montpellier II et Paris VII
$W$-fibrée sur $e$), donc $\emptyset \longrightarrow Y \times Z$ est dans $W$. Donc aussi la projection
$\begin{tikzcd}[column sep=small, row sep=small, baseline=(\tikzcdmatrixname-1-1.base)]
Y \times Z \arrow[rr] & & Y. \\
& \emptyset \arrow[ul] \arrow[ur] &
\end{tikzcd}$
Comme cette projection est une $W$-fibration, il s'ensuit par W(4b$''$) que ses fibres sont $W$-asphériques. Comme $Y \neq \emptyset$, il y a une telle fibre, elle est $\simeq Z$. Donc $Z$ est $W$-asphérique. Donc on trouve que \textit{toute} catégorie $Z$ dans Cat est $W$-asphérique. Donc toute flèche $Z \longrightarrow Z'$ dans Cat est dans $W$, q.e.d.

Supposons maintenant $\pi_0(f)$ surjective, mais non injective. Il y aurait une fibre de $\pi_0(X) \longrightarrow \pi_0(Y) = I$ qui serait de cardinal $> 1$. En prenant le $f_i$ correspondant, on trouve $X_i \longrightarrow Y_i$ avec $Y_i$ 0-connexe et $\card \pi_0(X_i) \geq 2$. Donc on est ramené au cas
\[
f : X \longrightarrow Y \qquad\qquad \left\{
\begin{array}{l}
f \in W \\[1ex]
Y \text{ 0-connexe} \\[1ex]
\card \pi_0(X) \geq 2.
\end{array}
\right.
\]
Considérons la factorisation canonique
\[
X \xrightarrow{i_f} \bar{X} \xrightarrow{\bar{f}} Y,
\]
avec $i_f \in W$, donc $\bar{f} \in W$. Donc par W(4b$''$) les fibres de $\bar{f}$ sont $W$-asphériques. Mais on voit tout de suite, en décomposant $\bar{X}$ suivant les composantes connexes $X_i$ de $X$, que ces fibres ont un $\pi_0$ de cardinalité $\geq \card \pi_0(X) \geq 2$. Donc on est ramené au cas où $Y = e$, i.e.
\[
X \ W\text{-asphérique}, \qquad \card \pi_0(X) \geq 2.
\]
Considérons $X \times X$, qui est aussi $W$-asphérique, et l'application diagonale\msp{page 22}
\[
d : X \longrightarrow X \times X = Y,
\]
qui est donc dans $W$. Mais $\pi_0(d)$ s'identifie à l'application diagonale $\pi_0(X) \longrightarrow \pi_0(X) \times \pi_0(X)$, qui n'est pas surjective puisque $\card \pi_0(X) \geq 2$. Donc pour la première partie, $W = \Fl(\Cat)$.

Donc sous forme plus positive :

\textbf{Corollaire.} \textit{Soit} $W \subseteq \Fl(\Cat)$ \textit{qui satisfait à} W(1,2,3,4b$''$)\textit{, et telle que} $W \neq \Fl(\Cat)$. \textit{Alors $W$ satisfait l'axiome de connexité} W(5)\textit{, et par suite (th. 1, page 12)} $W \subseteq W_\infty$.

\end{document}
